\begin{abstract}
Can a behaviour-based audit distinguish an instructed cartel from honest
language-model bidders? We test pairwise net influence in a repeated
procurement auction, comparing public-bid access with access to private
costs. The sole confirmatory test distinguished strong-pair cartel markets
from matched honest controls: AUROC $0.865$, descriptive interval
$[0.734,0.968]$, and $p=0.00002$ across $20$ blocks.
% source: results/analysis/final_v1/results.json primary; all_cells[206].auroc
Same-access residual correlation achieved $0.9975$.
% source: results/analysis/final_v1/results.json all_cells[202].auroc
Net influence therefore shows no general advantage over correlation.
In descriptive secondary comparisons, espionage was the only condition
whose paired net-minus-correlation interval excluded zero in favour of net
influence. Market-level flow labels remained largely weak or unclear.
The Qwen judge returned invalid JSON for $322/600$ judgements;
its matched comparisons cover only $1$--$8$ blocks.
% source: results/analysis/final_v1/results.json judgments grouped by model/status/failure; judge_comparisons[*].judge_auroc.n_blocks
Exploratory probes did not resolve an average linear bid response to partner
cost. Their leak test has an anti-conservative bid-level permutation null
because bids share regenerated messages. These results support screening
for an assigned treatment in this game, not reliable identification of
agreement, intent, or economic harm.
\end{abstract}
